%
% $Log: goodproblem.tex,v $
% Revision 1.7  2002/10/21 15:26:44  mjm
% to ohio, revised
%
% Revision 1.6  2002/01/08 18:58:40  mjm
% changed fonts
%
% Revision 1.5  2001/08/15 18:07:03  mjm
% APPM preprint 466
%
% Revision 1.4  1999/08/09 19:55:27  mjm
% beta release
%
% Revision 1.3  1999/04/17 21:33:11  mjm
% Kristian comments
%
% Revision 1.2  1999/04/06 19:22:02  mjm
% Ellens comments
%
% Revision 1.1  1999/03/17 00:21:00  mjm
% good problems
%
%

\documentclass{article}

\usepackage{amsmath,amsxtra,amssymb,epsfig,eepic} 
%\pagesize{usletter} 
\newlength{\originalbase}
\setlength{\originalbase}{\baselineskip}
\newcommand{\spacing}[1]{\setlength{\baselineskip}{#1\originalbase}}
%% the spacing command allows doublespacing etc.  Since there is
%% already some extra space built in, we don't call it with argument
%% 2 for doublespacing, but instead by the following chart:
%%      10pt type       \spacing{1.67}
%%      11pt type       \spacing{1.62}
%%      12pt type       \spacing{1.66}
%% To return to normal spacing, use     \spacing{1}
 
 
%%%% stretches the rows in arrays. 1.67 gives doublespace%%%%%%%%%
\renewcommand{\arraystretch}{1.1}
 
%%%%%%%%%%%%%%%%%%%%%%%%%%%%%%%%%%%%%%%%%%%%%%%%%%%%%%%%%%%%%%%%%%%%%%%%%%%%
%%%%% puts section headings and page number at the top. not required %%%%%%
%\pagestyle{headings}
%% if you do NOT use \pagestyle{headings} set to 0in
\setlength{\headheight}{0in}
\setlength{\headsep}{0in}
%%%%%%%%%%%%%%%%%%%%%%%%%%%%%%%%%%%%%%%%%%%%%%%%%%%%%%%%%%%%%%%%%%%%%%%%%%%%
 


%%%%%%%%% SET UP MARGINS %%%%%%%%%%%%%%%%%%%%%%%%%%%%%%%%%%%%%%%%%%%%%%%%%%%
%% left margin 1.25in (1in is automatic) right margin 1in
\setlength{\oddsidemargin}{0.25in}
%% evenside only counts for two sided printing
%% this leaves 1.25in on the left
\setlength{\evensidemargin}{0.25in}
\setlength{\textwidth}{\paperwidth}
        \addtolength{\textwidth}{-2.5in}
%% top margin 1in (automatic)
\setlength{\topmargin}{0in}
%%  bottom margin 1in
\setlength{\textheight}{\paperheight}
        \addtolength{\textheight}{-2in}
        \addtolength{\textheight}{-\headheight}
        \addtolength{\textheight}{-\headsep}
%% page numbers at bottom are 0.5in down into the 1in bottom margin
%% and so 0.5in from the bottom
\setlength{\footskip}{0.5in}
%% marginpars should only be used in draft copies
\setlength{\marginparsep}{1em}
\setlength{\marginparwidth}{0.6in}
%%%%%%%%%%%%%%%%%%%%%%%%%%%%%%%%%%%%%%%%%%%%%%%%%%%%%%%%%%%%%%%%%%%%%%%%%%%%
 
\newenvironment{examplelist}{\begin{enumerate}}{\end{enumerate}}

\newenvironment{example}{\item }{}

\newenvironment{badexample}{
%  \fontencoding{T1}\fontfamily{cmfr}\selectfont 
	\noindent \textbf{Bad:}\sl }{}

\newenvironment{goodexample}{
%  \fontencoding{OT1}\fontfamily{pss}\selectfont 
	\noindent \textbf{Good:}\sf }{}

%%%%%%%%%%%%%%%%%%%%%%%%%%%%%%%%%%%%%%%%%%%%%%%%%%%%%%%%%%%%%%%%%%%%%%%%%%%%
\newcommand{\honame}{{handout name}}
\newcommand{\handout}[1]{
	\newpage
	\renewcommand{\honame}{{#1}}
	{\noindent {\Large {\bf \honame}} 
		\hfill {\tiny Good Problems: \today}\\
	}
}
\newcommand{\secondpage}{
	\newpage
	{\noindent { {\bf \honame}, page 2.} 
		\hfill {\tiny Good Problems: \today}\\
	}
}
%%%%%%%%%%%%%%%%%%%%%%%%%%%%%%%%%%%%%%%%%%%%%%%%%%%%%%%%%%%%%%%%%%%%%%%%%%%%

\title{Good Problems:
teaching mathematical writing}
%\author{}
\date{Revised \today}

%%%%%%%%%%%%%%%%%%%%%%%%%%%%%%%%%%%%%%%%%%%%%%%%%%%%%%%%%%%%%%%%%%%%%%%%%
%%%%%%%%%%%%%%%%%%%%%%%%%%%%%%%%%%%%%%%%%%%%%%%%%%%%%%%%%%%%%%%%%%%%%%%%%
%%%%%%%%%%%%%%%%%%%%%%%%%%%%%%%%%%%%%%%%%%%%%%%%%%%%%%%%%%%%%%%%%%%%%%%%%
\begin{document}

\spacing{1.67}

\maketitle

\vfill
\subsection*{Abstract:}

Many students, especially those in science and engineering, are not
able to effectively communicate mathematical ideas in
writing. Furthermore, they generally believe that this ability is not
important. Employers, however, consider this to be a crucial
skill. Instructors would have to side with the employers, but usually
do not attempt to {\em teach} the students how to write, or that
writing is important.

We present here a plan to make writing coherent mathematics a regular
part of the students' experience.  This will both improve their
writing and teach them that writing is valued in mathematics.  We
provide materials to teach a set of important writing skills, so the
students will not have to guess how to write. This method is designed
to be as painless as possible, for the instructors as well as for the
students.

\vfill
\spacing{1.0}
\subsection*{Credits:}

The Good Problems are maintained and updated by Martin J. Mohlenkamp;
Department of Mathematics, 
	Ohio University, 321 Morton Hall, Athens OH 45701; 
	\texttt{mjm@math.ohiou.edu}.

They were originally developed in the Spring of 1999 in the Teaching
and Learning Seminar in the Department of Applied Mathematics at the
University of Colorado.% An early version 
\begin{description}
\item[Good Problems: teaching mathematical writing.]
Danielle Bundy, Ellen Gibney, Jenny McColl, Martin J. Mohlenkamp,
Kristian Sandberg, Brian Silverstein, Peter Staab, and
Matthew Tearle.
{\em University of Colorado APPM preprint \#466,} August 15, 2001;\\
\texttt{http://amath.colorado.edu/activities/preprints/}
\end{description}


%\marginpar{?copyright}

%\vfill
%Unresolved:
%\begin{itemize}
%How can the instructor do quality control of the TAs? (Perhaps a more
%general problem.)
%	\item
%We should mention any problems that we can think of that might occur in
%implementing this method, and give advice on overcoming them.
%\end{itemize}


\spacing{1.3}
\input{method}

\spacing{1.5}
%
% $Log: contents.tex,v $
% Revision 1.6  2002/10/21 15:26:44  mjm
% to ohio, revised
%
% Revision 1.5  2001/08/15 18:07:03  mjm
% APPM preprint 466
%
% Revision 1.4  1999/08/09 19:55:27  mjm
% beta release
%
% Revision 1.3  1999/04/06 19:29:00  mjm
% Matts comments
%
% Revision 1.2  1999/04/06 19:22:02  mjm
% Ellens comments
%
%

%%%%%%%%%%%%%%%%%%%%%%%%%%%%%%%%%%%%%%%%%%%%%%%%%%%%%%%%%%%%%%%%%%%%%%%%%
\handout{How to Use This Packet}

This packet contains all the materials necessary to implement the good
problems program. The items included are:
\begin{description}
	\item[The Instructor's Guide.] This guide is for the main
instructor, meaning the person in charge of organizing the course. It lists
those tasks that the 
instructor must perform to implement this scheme, indicates when they must be
performed, and estimates how much time they are expected to take.
If there is more than one instructor, then several of these tasks can be
shared. 
	\item[The Teaching Assistant's Guide.]
This guide is for the teaching assistant, who runs the recitations.
It lists those tasks that the teaching assistant must perform, indicates
when they must be performed, and estimates the time required.
If the course is taught without recitations, then these tasks
fall on the instructor.
	\item[The Student's Guide.] This guide explains the method to the
students, and tells them what is expected from them.
When this program is implemented, this guide could be one side of a sheet
of paper and the schedule of good problems for the course photocopied onto
the back. 
	\item[The Handouts.] The main material for this method is a set of
six handouts. 
Each handout is two pages long, so that it can be copied double-sided onto
a single sheet of paper. 
The handout consists of a description of a particular writing skill,
motivation for the importance of this skill,
whatever explicit rules can be given,
and examples of good and bad presentation emphasizing that skill.
 It is hoped that
the students will find these handouts useful enough that they will keep
them for future reference.

The handouts could be used in any order, but we suggest the following sequence:
\begin{enumerate}
\item Laying out the Problem.
\item Flow.
\item Mathematical Symbols.
\item Logical Connectives.
\item Graphs.
\item Introductions and Conclusions.
\end{enumerate}

\end{description}



%%% begin guide handouts



\spacing{1.4}
%
% $Log: instguide.tex,v $
% Revision 1.6  2002/10/21 15:26:44  mjm
% to ohio, revised
%
% Revision 1.5  2001/08/15 18:07:03  mjm
% APPM preprint 466
%
% Revision 1.4  1999/08/09 19:55:27  mjm
% beta release
%
% Revision 1.3  1999/04/06 19:29:00  mjm
% Matts comments
%
% Revision 1.2  1999/04/06 19:22:02  mjm
% Ellens comments
%
%

%%%%%%%%%%%%%%%%%%%%%%%%%%%%%%%%%%%%%%%%%%%%%%%%%%%%%%%%%%%%%%%%%%%%%%%%%
\handout{Instructor's Guide to Good Problems}

 This guide is for the main instructor, meaning the person in charge of
organizing the course. It lists those tasks that the instructor must
perform to implement this scheme, indicates when they must be performed,
and estimates how much time they are expected to take.
If there is more than one instructor, then several of these tasks can be
shared.

The first time you use this method, there are extra tasks, which
we list first.
\begin{enumerate}
	\item Read this packet, learn this material, make modifications,
and decide that it is a good thing to do. This should be performed well
before the term starts, to allow the ideas to settle. Time required: 3
hours.
	\item Convince the other people involved with the course to go
along with this plan. This task should be completed a month before the
term starts. Time required: 1 hour.
	\item Deal with unanticipated problems. Time required: 2 hours.
\end{enumerate}

Each time the method is used, there are the following tasks.
\begin{enumerate}
\item 
Choose the good problems, schedule them, and decide on the time
schedule for each new skill. The problems should not be harder than
the regular homework problems, but should be substantive enough that
there is something to write about. They must also be chosen to be
relevant to the new skill learned that week. (If the handout is on
graphing, the problem must require a graph.) Since there are only six
handouts, there is flexibility in
scheduling certain skills for certain sections, or avoiding busy
(test) weeks.  You might consider making the good problem on the
previous week's material, since a student will not be able to write
about new material that they do not yet understand.  This task should
be performed when the regular homework problems are chosen. Time
required: 1 hour.  \item Provide initial training for teaching
assistants. You must introduce this method, distribute materials and
discuss its implementation with the teaching assistants. This should
occur one week before the term starts. Time required: 2 hours.
\item Explain the method to the students. The concept should be
introduced and the student's guide distributed in the first week of
classes. The skills handouts could be distributed all at once at the
beginning or as they are used. Time required: 10 minutes of class
time.  \item Provide feedback on the handouts. Any deficiencies in the
method or the handouts should be noted. At the end of the term
comments from the instructors, teaching assistants, and students
should be collected and the materials corrected. Time required: 1
hour.
\end{enumerate}

{\bf Remark:}
Grade the good problem either as a significant portion of the homework or as a
quiz. A balance must be struck so the students will put effort into the
good problem without neglecting the others. 






\pagestyle{empty} % in effect for instguide


\spacing{1.5}
%
%$Log: taguide.tex,v $
%Revision 1.6  2002/10/21 15:26:44  mjm
%to ohio, revised
%
%Revision 1.5  2001/08/15 18:07:03  mjm
%APPM preprint 466
%
%Revision 1.4  1999/08/09 19:55:27  mjm
%beta release
%
%Revision 1.3  1999/04/06 19:29:00  mjm
%Matts comments
%
%Revision 1.2  1999/04/06 19:22:02  mjm
%Ellens comments
%
%


%%%%%%%%%%%%%%%%%%%%%%%%%%%%%%%%%%%%%%%%%%%%%%%%%%%%%%%%%%%%%%%%%%%%%%%%%
%\handout{Teaching Assistant's Guide to Good Problems}
\handout{TA's Guide to Good Problems}

This guide is for the teaching assistant, who runs the recitations.
It lists those tasks that the teaching assistant must perform, indicates
when they must be performed, and estimates the time required.
If the course is taught without recitations, then these tasks
fall on the instructor.

The first time you use this method, there are extra
tasks, which we list first.
\begin{enumerate}
	\item Read and understand the material in this packet. This
should be performed in the week before the term starts.
Time required: 2 hours.
	\item Attend training. The instructor should provide some training
during the week before the term starts. Time required: 1 hour.
	\item Deal with unanticipated problems. Time required: 2 hours.
\end{enumerate}

Each time the method is used there are the following tasks.
\begin{enumerate}
	\item Prepare for the new skill to be taught that week. It may also
be useful to review the old skills. This is part of the weekly preparation
time. Time required: 5 minutes per week.
	\item Present difficult skills in the recitation. Sometimes
students will have particular difficulty on some specific skill. The
teaching assistant will need to re-explain this skill and perhaps
demonstrate it on a relevant problem. It is important to note that the
teaching assistant is not expected to be the main teacher of these
skills. The handouts should provide all the information the students
need. Time required: 5 minutes recitation time per week, on average.
	\item Grade the good problems. It seems fair to
grade two less regular homework problems to compensate for the good
problems. After some initial trauma, it is expected that the good problems
will be more pleasant to grade than ordinary problems. There should be no
net gain in grading time as a result of this.
	\item Provide feedback on the handouts. Any deficiencies in the
method or the handouts should be noted. At the end of the term these
comments should be collected by the instructor and the materials
corrected. Time required: 1 hour. 
\end{enumerate}


\spacing{1.4}
%
% $Log: studguide.tex,v $
% Revision 1.5  2002/10/21 15:26:44  mjm
% to ohio, revised
%
% Revision 1.4  2001/08/15 18:07:03  mjm
% APPM preprint 466
%
% Revision 1.3  1999/08/09 19:55:27  mjm
% beta release
%
% Revision 1.2  1999/04/06 19:22:02  mjm
% Ellens comments
%
%

%%%%%%%%%%%%%%%%%%%%%%%%%%%%%%%%%%%%%%%%%%%%%%%%%%%%%%%%%%%%%%%%%%%%%%%%%
\handout{Student's Guide to Good Problems}

Many students believe that writing has nothing to
do with mathematics. {\bf This is false.}
In any field, it is important to be able to communicate your ideas and
results to others. Exactly how you communicate them varies from field to
field, but in nearly all fields it is through writing. Mathematical
writing, however, has its own particular style. The emphasis is on clarity
and precision, and not on the clever turn of phrase.

The ``Good Problems'' program is designed to teach you how to write
about mathematics coherently. It provides you a set of specific
writing skills and gives you enough practice through regular
assignments so that writing well will become a habit. It is also
designed to be as painless as possible.

Throughout the term, some homework problems will be designated 
``good problems''. You  need to write the solution to this problem
carefully, in good ``presentation'' format. It will be graded partly for
correctness, but with emphasis on presentation.

During the term you will receive six handouts, each addressing a
particular writing skill. The good problem is graded only on those
skills already covered, with greater emphasis on the newest skill. By
the end of the term you will have acquired the basic skills of
mathematical writing. The handouts are:
\begin{itemize}
\item Laying out the Problem.
\item Flow.
\item Mathematical Symbols.
\item Logical Connectives.
\item Graphs.
\item Introductions and Conclusions.
\end{itemize}
Each handout consists of a description of a particular writing skill,
motivation for the importance of this skill,
whatever explicit rules can be given,
and examples of good and bad presentation emphasizing that skill.
We use a consistent format for the examples.

\smallskip
\begin{badexample}
An example of poor writing looks like this.
\end{badexample}

\smallskip
\begin{goodexample}
An example of good writing looks like this.\\
\fbox{\rm A comment within an example looks like this.}
\end{goodexample}

\smallskip
If you already have these basic writing skills, then it should only take you
ten extra minutes to write up a good problem. If you do not have
these skills, then you will have to spend extra effort in the short term to
acquire them. In the long term, however, these skills will save you much
time and effort. Besides making future writing assignments easier, these
skills will help you with problem-solving (by encouraging organization and
logic) and with reading mathematics. These skills should be carried into
future mathematics courses, and become a useful life skill.




%%% begin main handouts

\spacing{1.0}
%\Large

%
% $Log: layout.tex,v $
% Revision 1.4  2002/10/21 15:26:44  mjm
% to ohio, revised
%
% Revision 1.3  2001/08/15 18:07:03  mjm
% APPM preprint 466
%
% Revision 1.2  1999/08/09 19:55:27  mjm
% beta release
%
% Revision 1.1  1999/07/12 22:04:36  mjm
% Initial revision
%
%
%

%%%%%%%%%%%%%%%%%%%%%%%%%%%%%%%%%%%%%%%%%%%%%%%%%%%%%%%%%%%%%%%%%%%%%%%%%
%\subsection*{Laying out the Problem}

\handout{Laying out the Problem}

Presentation is very important in mathematics, just as it is in other
fields. The main reason for this is clarity. Good presentation allows you
to  communicate more clearly the content of your work.
A secondary reason is the appearance of competence. Careful presentation
makes the reader think you were also careful with the content of your
work. Although good presentation is rarely successful at bluffing past poor
content, poor presentation can easily ruin good content.

The first thing the reader evaluates is the overall visual layout of the
problem. Exactly how polished this should be will depend on the purpose of
your writing. For our purposes, we have the following rules:
\begin{itemize}
	\item
Put the good problem on its own sheet of paper, separate from any
other problems. Regular quality paper is fine, 
but there should be no ragged edges or tears.
	\item
If you need more than one page, staple (no folded corners!) them together
and put your name and the page number on each page.
	\item
Leave margins on all sides.
	\item
Print neatly or type. Do not switch colors or from pen to pencil in the
middle of the problem. (You can use different colors to highlight if you wish.)
	\item
If you had to cross out material or erased a lot and left smudges, rewrite
the problem. (It is a good habit to solve the problem first on scrap paper
and then copy it neatly.)
\end{itemize}
%
The reader next needs to know what it is they are reading. On top of the first
page, put:
\begin{enumerate}
	\item
your full name,
	\item
the course and section or recitation number, and
	\item
the assignment number and/or the date the assignment is due.
\end{enumerate}
%
Next is the format for the problem itself:
\begin{itemize}
	\item
Label the problem with the chapter, section, and problem number.
	\item
Write out the entire question, including any instructions. If the question
refers to another problem, include the relevant information from that
problem. The goal is to make your work as self-contained as possible, so
the reader does not need to look anything up. 
	\item
Do the problem in some logical order. Do not do the problem in several
disjoint pieces connected by arrows.
\end{itemize}

All these rules may seem picky. Once you have learned this way to lay out
your work, you will understand the principles behind these rules. Then you
will know when to change the rules.

{\bf Acknowledgements:} If you had help with this problem, or in any
way include work that is not your own, then give proper credit. Not
only is this the nice thing to do, but also it will protect you from
allegations of cheating.

Good presentation is important in mathematics. It lets you communicate your
work more clearly and lends it credibility.


%%%%%%%%%%%%%%%%%%%%%%%%%%%%%%%%%%%%%%%%%%%%%%%%%%%%%%%%%%%%%%%%%%%%%%%%%
\secondpage

\subsection*{Example:}
%{\bf Example:}

\begin{goodexample}

\begin{flushright}
Full Name\\
Course and section or recitation number\\
Assignment number\\
Date\\
\end{flushright}

\noindent{\bf Section number, problem number.}
 Full text of the problem \ldots blah \ldots blah \ldots blah \ldots blah.

\vspace{\baselineskip}

In this problem we \ldots blah \ldots blah \ldots blah \ldots blah \ldots
blah \ldots blah \ldots blah \ldots blah. \\ 
\fbox{\rm There is a separate handout on introductions.}

\vspace{\baselineskip}

Since \ldots blah \ldots blah \ldots we know \ldots blah \ldots blah.\\
\fbox{\rm There is a separate handout on logical connectives.}

\begin{align*}
\implies & \mbox{blah}\\
\implies & \mbox{blah blah}\\
\therefore & \quad \mbox{blah!}\\
& \fbox{\rm There is a separate handout on mathematical symbols.}
\end{align*}
From this we can conclude \ldots blah \ldots blah \ldots blah \ldots blah
and so $x=3\pi\ldots$ \ldots blah. \\
\noindent\fbox{\parbox{\textwidth}{\rm
There is a separate handout on incorporating mathematics into sentences and other issues of `flow'.}}

\vspace{\baselineskip}

By graphing our solution, we can see that \ldots blah.
\begin{center}
\setlength{\unitlength}{0.5mm}
\begin{picture}(100,100)(0,0)
	\put(0,0){\line(1,0){100}}
	\put(0,0){\line(0,1){100}}
	\put(100,0){\line(0,1){100}}
	\put(0,100){\line(1,0){100}}
	\put(8,60){\mbox{\rm There is a separate handout}}
	\put(20,40){\mbox{\rm on graphing.}}
\end{picture}
\end{center}


\vspace{\baselineskip}

By using \ldots blah \ldots blah \ldots we were able to solve this
problem \ldots blah \ldots blah \ldots blah \ldots blah.\\ 
\fbox{\rm There is a separate handout on conclusions.}

\vspace{\baselineskip}
{\bf Acknowledgements:} I worked with so-and-so on this problem and also had 
help from so-and-so.

\end{goodexample}



%
% $Log: sentences.tex,v $
% Revision 1.4  2002/10/21 15:26:44  mjm
% to ohio, revised
%
% Revision 1.3  2001/08/15 18:07:03  mjm
% APPM preprint 466
%
% Revision 1.2  1999/08/09 19:55:27  mjm
% beta release
%
% Revision 1.1  1999/06/10 23:34:48  mjm
% Initial revision
%
%

\handout{Flow}
       
Written mathematics must be readable. This may seem trivial, but it is
an important point. You should be able to read your work aloud to a
classmate and have them understand your solution. If you need to add
any explanations, these should be included in your written work.
The most common mistake is to write mathematics without using enough
words. All writing, even mathematics, should consist of complete sentences.
These should explain the problem by providing both the method and 
justification for each  step of the solution.


\begin{flushleft}
{\bf{Why are sentences important in mathematics?}}
\end{flushleft}

  Although sometimes it seems hard to read textbooks, it would be
much harder to understand if they only had equations and no sentences.
The situation is similar in lecture: if the professor just listed
formulas on the chalkboard without talking about them, leaving you to
figure out what was being done in each step, how much could you
understand from the lecture?  Neither of these would be a good way for
most students to learn, since sentences are necessary to explain the
mathematics.

\begin{flushleft}
  \bf{Why should students use sentences in a Mathematics class?}
\end{flushleft}

In a Mathematics class, you should explain homework solutions using complete
sentences.  That means linking together thoughts with words and
embedding equations into sentences.  Going through the extra work to
do this will benefit you in several ways:

\begin{itemize}
\item Writing down your thoughts and organizing them into complete
  sentences will help you to understand the method of solution better.
  
\item When you look back on homework to study for a test, or later on
  in another class, you will understand what you
  were doing on each problem and the mathematics behind it.
  
\item Other people (teacher, classmates, grader,...) will understand
  what you are doing at each step, and why you are doing it.  This
  way, you won't lose points for skipping steps or solving the problem
  in an unusual way.
  
\item Communicating your work will be essential in whatever field you choose.
  Even though the fields are stereotypically weak on
  writing, engineers and scientists spend a surprising amount of time
  writing reports and giving oral presentations.
\end{itemize}

%%%%%%%%%%%%%%%%%%%%%%%%%%%%%%%%%%%%%%%%%%%%%%%%%%%%%%%%%%%%%%%%%%%%%%%%%
\subsubsection*{The Royal ``We''.}

It is customary to write mathematics using ``we'' instead of ``I'' and
``you''. Think of yourself as the tour guide, showing the sights to
your reader. You and the reader together make ``we''.


\subsubsection*{Your Audience}
Keep in mind who your audience is. For the Good Problems, pretend that
you are writing for a classmate who has not seen this problem
before. Write with enough detail so that they can follow your
explanations.  A good way to test if you have written enough is to read
your work aloud to another person. If you need to add any words to
make it make sense, then these words should be included in your
written work.



%%%%%%%%%%%%%%%%%%%%%%%%%%%%%%%%%%%%%%%%%%%%%%%%%%%%%%%%%%%%%%%%%%%%%%%%%
\secondpage
\subsection*{{Examples:}}

\begin{itemize}

%\noindent{\bf How to put an equation into a sentence:}
%\subsubsection*{How to put an equation into a sentence:}
%
%\begin{goodexample}
%    The derivative of the curve $y= 3\sqrt{x}$ is $\frac{dy}{dx} =
%    \frac{3}{2 \sqrt{x}}$.
%\end{goodexample}  
%    
%\begin{badexample}
%  The derivative $\frac{dy}{dx}$ is $\frac{3}{2 \sqrt{x}}$ for the
%  curve $y$ which is $3 \sqrt{x}$.
%\end{badexample}
%    
%\begin{badexample}
%  $\frac{dy}{dx}$ is the derivative of $3 \sqrt{x}$ which is
%  $\frac{3}{2 \sqrt{x}}$.
%\end{badexample}



%\bigskip
%%%%%%%%%%%%%%%%%%%%%%%%%%%%%%%%%%%%%%%%%%%%%%%%%%%%%%%%%%%%%%%%%%%%%%%%%
%\subsubsection*{Using sentences incorrectly:}

\item  
Find an equation for the tangent to the curve $y = x + \frac{2}{x} $
at the point $(1,3)$.
 
%\begin{badexample}
%
%  \begin{displaymath}
%    \frac{dy}{dx} = 1 - \frac{2}{x^2}
%  \end{displaymath}
%                                %
%  so the equation for the derivative is
  %
%  \begin{displaymath}
%    \frac{dy}{dx}= -1 = m
%  \end{displaymath}
%                                %
%  so the slope at $(x,y)$ can be found by $y-3 = -1 (x-1)$ so $y=-x +
%  4$ is the point-slope equation.
%\end{badexample}
  

\begin{goodexample}
  
We first check that $(1,3)$ is a point on the curve by plugging these
values in: $3=1+2/1$. 
  The derivative of the curve $y = x + \frac{2}{x}$ is
%  \begin{displaymath}
\begin{equation}
    \frac{dy}{dx} = 1 - \frac{2}{x^2}.
%\label{eqn:refexample}
\end{equation}
%  \end{displaymath}
  A line tangent to the curve at the point $(1,3)$ will have slope
  \begin{displaymath}
    \left. \frac{dy}{dx} \right| _{x=1} = 1 - \frac{2}{(-1)^2} = -1.
  \end{displaymath}
  Using the point-slope formula with $m = -1$, $x_0 = 1$, and $y_0 = 3$ gives
  the formula for the line $y - 3 = -1(x-1)$.  Solving for y and
  simplifying gives
  \begin{displaymath}
    y = -x + 4.
  \end{displaymath}
  This is the equation of the line tangent to the curve at that point.
\end{goodexample}  


%%%%%%%%%%%%%%%%%%%%%%%%%%%%%%%%%%%%%%%%%%%%%%%%%%%%%%%%%%%%%%%%%%%%%%%%%
%\subsubsection*{Vague and uninformative sentences}
    
\item
\begin{badexample}
      We use calculus to find that $y = 3x^2+1$ has a slope of $3$ at
      $x = 1/2$.
\end{badexample}

\begin{goodexample}
    To find the slope of the curve $y = 3x^2+1$ at the point $x = 1/2$
    we find $\frac{dy}{dx}$ evaluated at $x = 1/2$. We compute
    \begin{eqnarray*}
      \frac{dy}{dx} &  = &  6x,  \\  
      \Rightarrow \left. \frac{dy}{dx} \right| _{x=1/2} & = & 6 \left( \frac{1}{2}
      \right) = 3. 
    \end{eqnarray*}
\end{goodexample}

\item
                 %not sure if this is a good example of
                 %not, seems a little contrived 
\begin{badexample}
    \begin{eqnarray*}
      \frac{d}{dx}\left( (x^2 + 1) (x^3 +3) \right) &  = & 
      (x^2+1)3x^2 +(x^3+3)2x \\ 
      &  = & 5x^4+3x^2+6x
    \end{eqnarray*}
    is the derivative.
\end{badexample}
    
\begin{goodexample}
    To take the derivative of a product of $2$ functions, we use the
    product rule, $(fg)'=f'g+fg'$.
In our case we have
    \begin{eqnarray*}
      \frac{d}{dx}\left( (x^2 + 1) (x^3 +3) \right) &  = & 
      \left(\frac{d}{dx} (x^2 + 1) \right) (x^3 +3)
	+(x^2 + 1) \left(\frac{d}{dx} (x^3 +3)\right)\\
	& =& (2x)(x^3+3) +(x^2 + 1)(3x^2)\\
      &  = & 5x^4+3x^2+6x.
    \end{eqnarray*}
\end{goodexample}

%%%%%%%%%%%%%%%%%%%%%%%%%%%%%%%%%%%%%%%%%%%%%%%%%%%%%%%%%%%%%%%%%%%%%%%%%
%\subsubsection*{Referring to previous equations and figures.}
\item
You may have noticed that one of the equations above has been labeled
equation number one by putting ``(1)'' at the right hand
margin. If you need to refer back to an equation or figure, label it
and then refer to it by its label. Do not draw arrows.

\begin{badexample}
Using the equation from before, the slope is $-1$. \fbox{\rm Which equation?}
\end{badexample}

\begin{goodexample}
Using (1), the derivative at $x=1$ is $-1$.
\end{goodexample}
\end{itemize}

%\marginpar{conclusion.}







%
% $Log: symbols.tex,v $
% Revision 1.5  2002/10/21 15:26:44  mjm
% to ohio, revised
%
% Revision 1.4  2002/01/08 18:54:53  mjm
% reworked table
%
% Revision 1.3  2001/08/15 18:07:03  mjm
% APPM preprint 466
%
% Revision 1.2  1999/08/09 19:55:27  mjm
% beta release
%
% Revision 1.1  1999/06/10 23:29:28  mjm
% Initial revision
%
%


%\begin{probdescript}{Using Mathematical Symbols}
\handout{Mathematical Symbols}


You will encounter many mathematical symbols during your math courses.
The table below provides you with a list  of the more common
symbols, how to read them, and notes on their meaning and usage.
The following page has a series of examples of these symbols in use.

\vfill

\renewcommand{\arraystretch}{1.2}
\noindent
%\begin{tabular}{|l|p{1.75in}|p{3.5in}|}\hline
\begin{tabular}{|p{0.8in}|p{1.6in}|p{3.15in}|}\hline
  {\bf Symbol} & {\bf How to read it} & {\bf Notes on meaning and usage} \\
  \hline
                                %
$ a=b $ &$a$ equals $b$& $a$ and $b$ have exactly the same value.  \\ \hline 
                                %
$a \approx b$ or \mbox{$a\cong b$} 
		& $a$ is approximately equal to $b$
			& Do not write $=$ when you mean $\approx$.\\ \hline
                                %
$ P\Rightarrow Q$ & $P$ implies $Q$ 
			&  If $P$ is true, then $Q$ is also true.\\ \hline
                                %
$ P\Leftarrow Q$ & $P$ is implied by $Q$ 
			&  If $Q$ is true, then $P$ is also true.\\ \hline
                                %
$ P\Leftrightarrow Q$ or $P$~iff~$Q$
		& $P$ is equivalent to $Q$ or $P$~if~and~only~if~$Q$
			& $P$ and $Q$ imply each other. \\ \hline 
                                %
$ (a,b)$ & the point $a$ $b$ 
		& A coordinate in $\mathbb{R}^2$.  \\ \hline
                                %
$ (a,b)$ & the open interval from $a$ to $b$ &
  	The values between $a$ and $b$, but
  	not including the endpoints.  \\ \hline
                                %
$ [a,b]$ & the closed interval from $a$ to $b$ &
  	The values between $a$ and $b$, including the endpoints.  \\ \hline 
                                %
$ (a,b]$ & The (half-open) interval from $a$ to $b$ excluding $a$, and
  	including $b$. 
		&The values between $a$ and $b$, excluding $a$,
		 and including $b$. Similar for $ [a,b)$. \\ \hline
                                %

$ \mathbb{R} $ or ${\bf R}$ & the real numbers 
		&It can also be used for the plane as $\mathbb{R}^2$, and
  		in higher dimensions. \\ \hline
                                %
$ \mathbb{C} $ or ${\bf C}$ & the complex numbers 
		& $\{a+bi: a,b\in \mathbb{R} \}$, where $i^2=-1$. \\ \hline
                                %
$ \mathbb{Z} $ or ${\bf Z}$ & the integers 
		& \ldots,$-2$,$-1$,0,1,2,3, \ldots. \\ \hline
                                %
$ \mathbb{N} $ or ${\bf N}$ & the natural numbers 
		&  $1,2,3,4, \ldots$. \\ \hline
                                %
$ a\in B$ & $a$ is an element of $B$
		& The variable $a$ lies in the set  (of values) $B$. \\ \hline
                                %
$a\notin B$ & $a$ is not an element of $B$ & \\ \hline
                                %
$A \cup B$ & $A$ union $B$
		& The set of all points that fall in $A$ or $B$.\\ \hline
                                %
$A \cap B$ & $A$ intersection $B$
		&The set of all points that fall in both $A$ and $B$.\\ \hline
                                %
$A \subset B$ & $A$ is a subset of $B$ or $A$~is~contained~in~$B$
	&Any element of $A$ is also an element of $B$.\\ \hline
				%
$ \forall x$ & for all $x$& Something is true for all (any) value of $x$
	(usually with a side condition like $ \forall x>0$).  \\ \hline
                                %
$\exists$   &  there exists  
		&Used in proofs and definitions as a shorthand. \\ \hline
                                %
$\exists!$  & there exists a unique 
		&Used in proofs and definitions as a shorthand.\\ \hline
                                %
$ f\circ g$ & $f$ composed with $g$ or $f$~of~$g$
			& Denotes $f(g(\cdot))$. \\ \hline
                                %
$n!$ & $n$ factorial & $n!=n(n-1)(n-2)\cdots \times 2 \times 1$. \\ \hline
                                %
$ \lfloor x \rfloor$ &the floor of $x$ 
		& The nearest integer $\le x$. \\ \hline
                                %
$ \lceil x \rceil$ &the ceiling of $x$ 
		& The nearest integer $\ge x$. \\ \hline
                                %
$f={\cal O}(g)$ or $f=O(g)$	& $f$ is big oh of $g$
		& $\lim_{x\rightarrow\infty}\sup_{y>x} |f(y)/g(y)|<\infty$. 
		Sometimes the limit is toward $0$ or another point.
					\\ \hline
				%
$f=o(g)$	& $f$ is little oh of $g$
		& $\lim_{x\rightarrow\infty}\sup_{y>x} |f(y)/g(y)|=0$.
					\\ \hline
				%
$x\rightarrow a^+$	& $x$ goes to $a$ from the right
		& $x$ is approaching $a$, but $x$ is always greater than $a$.
			Similar for $x\rightarrow a^-$.\\ \hline
\end{tabular}
\vfill
%%%%%%%%%%%%%%%%%%%%%%%%%%%%%%%%%%%%%%%%%%%%%%%%%%%%%%%%%%%%%%%%
\secondpage

%%%%%%%%%%%%%%%%%%%%%%%%%%%%%%%%%
\subsection*{The Trouble with $=$}
The most commonly used, and most commonly misused, symbol is
`$=$'. The `$=$' symbol means that the things on either side are
actually the same, just written a different way. The common misuse of
`$=$' is to mean 'do something'. For example, when asked to compute
$(3+5)/2$, some people will write:\\
\begin{badexample}
\[
3+5=8/2=4.
\]
\end{badexample}\\
This claims that $3+5=4$, which is false. We can fix this by carrying
the `$/2$' along, as in $(3+5)/2=8/2=4$. We could instead use the
'$\Rightarrow$' symbol, meaning 'implies', and turn it into a logical
statement: \\
\begin{goodexample}
\[
3+5=8\quad \Rightarrow \quad (3+5)/2=4.
\]
\end{goodexample}

%%%%%%%%%%%%%%%%%%%%%%%%%%%%%%%%%%%%%
\subsection*{To Symbol or not to Symbol?}

\begin{badexample}
$\lim_{x \rightarrow x_0} f(x)=L$ means that $\forall \epsilon >0$,
$\exists \delta >0$ s.t. $\forall x$,
\[
0 < |x -x_0|< \delta\quad  \Rightarrow\quad |f(x)-L|<\epsilon.
\]
\end{badexample}\\
Although this statement is correct mathematically, it is difficult to
read (unless you are well-versed in math-speak).  This example shows
that although you can write math in all symbols as a shortcut, often
it is clearer to use words. A compromise is often preferred.\\
\begin{goodexample}\\
{\bf The Formal Definition of Limit:}
Let $f(x)$ be defined on an open interval about $x_0$, except possibly
at $x_0$ itself. We say that $f(x)$ approaches the limit $L$ as $x$
approaches $x_0$, and we write
\[
\lim_{x \rightarrow x_0} f(x)=L
\]
if for every number $\epsilon>0$, there exists a corresponding number
$\delta>0$ such that for all $x$ we have
\[
0 < |x -x_0|< \delta \Longrightarrow |f(x)-L|<\epsilon.
\]
\end{goodexample}%\\
%The last statement says $0<|x -x_0|< \delta$ implies
%$|f(x)-L|<\epsilon$ is true.

%%%%%%%%%%%%%%%%%%%%%%%%%%
\subsection*{Other Examples}

The `$\Rightarrow$' symbol should be used even when doing simple algebra.\\
\begin{goodexample}
\[
(y-0)=2(x-1) \quad \Longrightarrow \quad y=2x-2
\]    
\end{goodexample}

You will be more comfortable with symbols, and better able to use them,
if you connect them with their spoken form and their meaning.\\
\begin{goodexample}
The mathematical notation $(f \circ g)(x)$ is read ``$f$ composed with
$g$ at the point $x$'' or ``$f$ of $g$ of $x$'' and means
\[
%        (f \circ g)(x) \quad \Leftrightarrow \quad 
f(g(x)).
\]
\end{goodexample}
    
%    \begin{example}{\bf Using $\Leftrightarrow$}\\
%      For sets of numbers on the real number line the following
%      two columns are equivalent:
%      \begin{eqnarray*}
%        x \in (a,b) &\Leftrightarrow & a < x < b \\
%        x \in [a,b] &\Leftrightarrow & a \leq x \leq b \\
%        x \in [a,b) &\Leftrightarrow & a \leq x < b \\
%        x \in (a,b] &\Leftrightarrow & a < x \leq b \\
%      \end{eqnarray*}
      
%    \end{example}    
                                % Example #3
    % example #4

%    \begin{example}
%      {\bf Theorem:  One-sided vs. Two-sided Limits}
%      
%      \nopagebreak
%      
%      A function $f(x)$ has a limit as $x$ approaches $c$ if and only
%      if it has left-hand and right-hand limits and they are equal:
%                                %
%      \begin{displaymath}
%        \begin{aligned}
%          \lim_{x\rightarrow c} f(x)  = L \quad \Leftrightarrow & \quad
%          \lim_{x\rightarrow c^{-}} f(x)= L \quad \text{and} \\
%          ~&\quad \lim_{x\rightarrow c^{+}} f(x)= L.
%        \end{aligned}
%      \end{displaymath}
%      
%      \begin{displaymath}
%           \lim_{x\rightarrow c} f(x)  = L \quad \Leftrightarrow  \quad
%          \text{both} \lim_{x\rightarrow c^{-}} f(x)= L \quad \text{and} 
%          \quad \lim_{x\rightarrow c^{+}} f(x)= L.
%       \end{displaymath}
%      
%      
%\noindent\fbox{\parbox{\textwidth}{\rm      
%        The above says that the statement to the left of the
%        $\Leftrightarrow$ is equivalent to the statement on the right
%        hand side.  
%        
%        Another way to look at this is that the statement on the left
%        implies the statement on the right, and conversely, the
%        statement on the right implies the statement on the left.
%}}
%    \end{example}
      
    % Example #5
  
    % Example #6
%
%
    
%  \end{examplelist}


%
% $Log: logicalcon.tex,v $
% Revision 1.4  2002/10/21 15:26:44  mjm
% to ohio, revised
%
% Revision 1.3  2001/08/15 18:07:03  mjm
% APPM preprint 466
%
% Revision 1.2  1999/08/09 19:57:38  mjm
% beta release
%
%

% Notes:
% 1. should we include inverse?

%%%%%%%%%%%%%%%%%%%%%%%%%%%%%%%%%%%%%%%%%%%%%%%%%%%%%%%%%%%%%%%%%%%%%%%%%
\handout{Logical Connectives}

%% reset equation counter
\setcounter{equation}{0}

Mathematics has its own language.  As with any language, effective
communication depends on logically connecting components.  Even the
simplest ``real" mathematical problems require at least a small amount
of reasoning, so it is very important that you develop a feeling for
{\em formal (mathematical) logic}.

Consider, for example, the two sentences ``There are 10 people waiting
for the bus" and ``The bus
is late."  What, if anything, is the logical connection between these two
sentences?  Does one {\em logically imply} the other?
Similarly, the two mathematical statements ``$r^2 + r - 2 = 0$" and
``$r=1$ or $r=-2$" need to be connected, otherwise they are merely two
random statements that convey no useful information.
  Warning: when
mathematicians talk about implication, it means that one thing {\bf must}
be true as a consequence of another; not that it can be true, or
might be true sometimes.  

Words and symbols that tie statements together logically are called
{\em logical connectives}.  They allow you to communicate the reasoning
that has led you to your conclusion.  Possibly the most important of
these is {\em implication} --- the idea that the next statement is a
logical consequence of the previous one.  This concept can be conveyed by
the use of words such as: therefore, hence, and so, thus, since, if $\dots$
then $\dots$, this implies, etc.  In the middle of mathematical calculations, we
can represent these by the implication symbol ($\Rightarrow$).
   For example
\begin{eqnarray}
x + 7y^2 = 3 & \Rightarrow & y = \pm \sqrt{\frac{3-x}{7}} \, ; \label{first} \\
x \in (0,\infty) & \Rightarrow & \cos(x) \in [-1,1]. \label{second}
\end{eqnarray}

\subsubsection*{Converse}
Note that ``statement A $\Rightarrow$ statement B" does {\bf not}
necessarily mean that the {\em logical converse} --- ``statement B
$\Rightarrow$ statement A" --- is also true.  Logical implication is a
matter of cause and effect; the logical converse is simply the reverse
cause-effect situation (which may not be true).
Consider the following, everyday example:
\begin{itemize}
\item ``I am running to class because I am late."
\end{itemize}
It should be clear that the logical converse of these is not true:
\begin{itemize}
\item ``I am late to class because I am running."
\end{itemize}

For examples (\ref{first}) and (\ref{second}) above:
\begin{eqnarray*}
x + 7y^2 = 3 & \Leftarrow & y = \pm \sqrt{\frac{3-x}{7}} \, ; \\
x \in (0,\infty) & \not\Leftarrow & \cos(x) \in [-1,1]. \quad
	\fbox{Since $x<0 \Rightarrow \cos(x) \in [-1,1]$ also.}
\end{eqnarray*}


\subsubsection*{Contrapositive}
We have seen that simply reversing a logical statement can lead to problems.
There is a way, however, to invert an implication so that the inverted
statement is also true.  This is known as the {\em contrapositive}.
Consider the statement ``$A \Rightarrow B$"; this means ``if $A$ is
true, then $B$ is true".  The contrapositive is ``if $B$ is {\bf false},
then $A$ must be false".  
%Recall our earlier example: a power outage causes
%everyone's TV to turn off.  If someone's TV is on ({\em i.e.} it is not
%true that everyone's TV is off), then there cannot be a power outage.

Here is an example of how these concepts fit together:
\begin{itemize}
\item If X is a cat, it is an animal.  \fbox{Accept as true}
\item If X is an animal, it is a cat.  \fbox{FALSE (converse)}
\item If X is not an animal, it is not a cat.  \fbox{TRUE (contrapositive)}
\item If X is not a cat, it is not an animal.\\
	  \fbox{FALSE. This is the
      contrapositive of the second statement, which is false.}
\end{itemize}

\secondpage


\subsubsection*{Equivalence}
When the implication works both ways, we say that the two statements
are {\em equivalent} and we may use the equivalence symbol
($\Leftrightarrow$); in words we may say ``A is equivalent to B" or
``A if and only if B".  If two statements are equivalent, we may use
any of the implication symbols ($\Rightarrow$, $\Leftarrow$, or
$\Leftrightarrow$).  Which connective we use depends on what we are
trying to show.  In (\ref{first}) above, if we are trying to obtain a
formula for $y$, we would probably just use ``$\Rightarrow$", even
though the stronger statement (``$\Leftrightarrow$") is also true.
If, however, we wanted to show that two statements about $x$ and $y$
were equivalent, we would write $x + 7y^2 = 3
\Leftrightarrow y = \pm \sqrt{\frac{3-x}{7}}$.


Careful logic is the heart and soul of mathematics; learn to reason with
watertight arguments and use logical connectives to explain your reasoning
processes.

\subsubsection*{Example:}

{\em What is the greatest amount of water that a right-cylindrical water tank
can hold if there is $100 m^2$ of material from which to construct it?}

\begin{goodexample}
Let the height of the tank be $h$ and the radius be $r$.  \underline{Then}
the volume of the tank is $V(h,r) = \pi hr^2$ and the surface area is
$S(h,r) = 2\pi r^2 + 2\pi rh = 2\pi r(r+h)$.  \underline{Since} we require
$V > 0$, we can assume that $r,h > 0$.  \underline{If} we have $100 m^2$
of material, \underline{then} $S = 100$, \underline{and so}
\begin{eqnarray*}
S & = & 2\pi r(r+h) = 100 \\
 & \Rightarrow & r + h = \frac{100}{2\pi r} \\
 & \Rightarrow & h = \frac{50}{\pi r} - r \\
 & \Rightarrow & V = \pi hr^2 = 50r - \pi r^3.
\end{eqnarray*}
To find the maximum volume, we look for $r$ such that $\frac{dV}{dr} = 0$.
Differentiating with respect to $r$ gives $\frac{dV}{dr} = 50 - 3\pi r^2$.
\underline{Hence},
\begin{eqnarray*}
\frac{dV}{dr} = 0 & \Leftrightarrow & 50 - 3\pi r^2 = 0 \\
    & \Leftrightarrow & r^2 = \frac{50}{3\pi} \\
    & \Leftrightarrow & r = \pm \sqrt{\frac{50}{3\pi}}.
\end{eqnarray*}
\underline{Since} $r>0$, we can use the second derivative test and find 
$\frac{d^2V}{dr^2} = -6\pi r <0$. 
\underline{This implies} that  $r = \sqrt{\frac{50}{3\pi}} \approx 2.3033$
maximizes $V$.
From above, the optimal volume is
\begin{eqnarray*}
V & = & 50r - \pi r^3 \\
  & = & \left( 50 - \pi \left( \frac{50}{3\pi} \right) \right) \sqrt{\frac{50}{3\pi}} \\
  & = & \left( \frac{100}{3} \right) \sqrt{\frac{50}{3\pi}} \approx 76.7765.
\end{eqnarray*}
\underline{Thus} we have found that the greatest amount of water such a tank
can hold is (approximately) $76.7765 m^3$.
\end{goodexample}




%
% $Log: graphs.tex,v $
% Revision 1.5  2002/10/21 15:26:44  mjm
% to ohio, revised
%
% Revision 1.4  2002/01/08 18:53:18  mjm
% made pdf compatible
%
% Revision 1.3  2001/08/15 18:07:03  mjm
% APPM preprint 466
%
% Revision 1.2  1999/08/09 19:55:27  mjm
% beta release
%
% Revision 1.1  1999/06/10 23:30:48  mjm
% Initial revision
%
%

%\begin{probdescript}{Drawing Graphs}
\handout{Graphs}

Graphs are an important visual aid in mathematics and it is important to
learn how to draw them properly.  For a graph to be useful it must display
the important and interesting information about a function and have all
relevant points clearly labeled.  

When drawing graphs, it is helpful to
ask yourself:
\begin{itemize}
\item Have I graphed what is being asked?  
\item Have I put a title on the  graph and labeled the axes?  
\item Is the scale I have chosen for my graph appropriate?
\item Have I illustrated the important features?
\item Have I included enough detail?
%  For example, labels on extreme
%  points, points of inflection, and $x$- and $y$-intercepts.
\end{itemize}

You can often determine how much detail to include in your graph from the
question.  For example, if you are finding the extreme points of a
function, then you will probably be expected to label these points on your
graph.  If you are only sketching the graph of a function then less detail
is required.  The main purpose of a sketch is to get some idea of the
general shape and behavior of the function.  In this case, it is not
usually necessary to label extreme points and intercepts.  However, you
should always remember to label your axes and put a title on your graph,
even if it is only a sketch.  Use a ruler if necessary to draw
straight lines. If 
you are drawing more than one curve on the same set of axes, take
particular care to make sure that each curve is clearly labeled.

%\end{probdescript}
%\clearpage

%\begin{probexamples}
%\begin{examplelist}
%\begin{example}

%\noindent{\bf Examples:}
\subsection*{Examples:}

  Graph the curve $y= x^3 - 4x$.

\begin{goodexample}

\setlength{\unitlength}{1cm}
\begin{center}
The curve $y=x^3-4x$.\\
\smallskip
\begin{picture}(10,8)(-5,-4)

% make axes

\put(0,-4){\vector(0,1){8}}
\put(-5,0){\vector(1,0){10}}
\thinlines
\multiput(-4,-0.1)(1,0){9}{\line(0,1){0.2}}

\multiput(-0.1,-3.5)(0,0.5){15}{\line(1,0){0.2}}

\put(5.2,0.2){$x$}
\put(0.2,3.8){$y$}

\put(-4.4,-0.4){$-4$}
\put(-3.4,-0.4){$-3$}
\put(-2.4,-0.4){$-2$}
\put(-1.4,-0.4){$-1$}

\put(0.9,-0.4){$1$}
\put(1.9,-0.4){$2$}
\put(2.9,-0.4){$3$}
\put(3.9,-0.4){$4$}

\put(-0.7,-2.6){$-5$}
\put(-0.4,2.4){$5$}
% add in the curve x^3+4x

%% FAKE but looks same

\qbezier(-2.6,-3.6)(-1.7,4.35)(0,0)
\qbezier(0,0)(1.7,-4.35)(2.6,3.6)

%\path(-2.6,-3.588)
%(-2.5,-2.8125)
%(-2.4,-2.112)
%(-2.3,-1.4835)
%(-2.2,-0.924)
%(-2.1,-0.4305)
%(-2,0)
%(-1.9,0.3705)
%(-1.8,0.684)
%(-1.7,0.9435)
%(-1.6,1.152)
%(-1.5,1.3125)
%(-1.4,1.428)
%(-1.3,1.5015)
%(-1.2,1.536)
%(-1.1,1.5345)
%(-1,1.5)
%(-0.9,1.4355)
%(-0.8,1.344)
%(-0.7,1.2285)
%(-0.6,1.092)
%(-0.5,0.9375)
%(-0.4,0.768)
%(-0.3,0.5865)
%(-0.2,0.396)
%(-0.1,0.1995)
%(0,0)
%(0.1,-0.1995)
%(0.2,-0.396)
%(0.3,-0.5865)
%(0.4,-0.768)
%(0.5,-0.9375)
%(0.6,-1.092)
%(0.7,-1.2285)
%(0.8,-1.344)
%(0.9,-1.4355)
%(1,-1.5)
%(1.1,-1.5345)
%(1.2,-1.536)
%(1.3,-1.5015)
%(1.4,-1.428)
%(1.5,-1.3125)
%(1.6,-1.152)
%(1.7,-0.9435)
%(1.8,-0.684)
%(1.9,-0.3705)
%(2,0)
%(2.1,0.4305)
%(2.2,0.924)
%(2.3,1.4835)
%(2.4,2.112)
%(2.5,2.8125)
%(2.6,3.588)

% add the labels

%\put(-2,0){\circle*{0.1}}
\put(-2.08,-0.08){\mbox{$\bullet$}}
\put(-3.2,0.3){$(-2,0)$}

%\put(-1.1,1.5396){\circle*{0.1}}
\put(-1.18,1.45){\mbox{$\bullet$}}
\put(-2.5,2){$(\frac{-2}{\sqrt{3}}, \frac{16}{3\sqrt{3}})$}


%\put(0,0){\circle*{0.1}}
\put(-0.08,-0.08){\mbox{$\bullet$}}
\put(0.1,0.5){$(0,0)$}

%\put(1.1,-1.5396){\circle*{0.1}}
\put(1.02,-1.62){\mbox{$\bullet$}}
\put(0.5,-2){$(\frac{2}{\sqrt{3}}, \frac{-16}{3\sqrt{3}})$}

%\put(2,0){\circle*{0.1}}
\put(1.92,-0.08){\mbox{$\bullet$}}
\put(2.4,0.5){$(2,0)$}



\end{picture}
\end{center}
\noindent\fbox{\parbox{\textwidth}{\rm The curve,
axes, extreme points, and $x$- and $y$-intercepts are clearly labeled.
These are the `points of interest' for this function.}}

\end{goodexample}

\secondpage

\begin{badexample}

\setlength{\unitlength}{0.7cm}
\begin{center}
\begin{picture}(10,8)(-5,-4)

% make axes

\put(0,-4){\vector(0,1){8}}
\put(-5,0){\vector(1,0){10}}
\thinlines
\multiput(-4,-0.1)(1,0){9}{\line(0,1){0.2}}

%\multiput(-0.1,-3.5)(0,0.5){15}{\line(1,0){0.2}}

\put(5.2,0.2){$x$}
\put(0.2,3.8){$y$}

\put(-4.4,-0.4){$-4$}
\put(-3.4,-0.4){$-3$}
\put(-2.4,-0.4){$-2$}
\put(-1.4,-0.4){$-1$}

\put(0.9,-0.4){$1$}
\put(1.9,-0.4){$2$}
\put(2.9,-0.4){$3$}
\put(3.9,-0.4){$4$}

% add in the curve x^3+4x

% FAKE 
\qbezier(-2.6,-3.6)(-1.7,4.35)(0,0)
\qbezier(0,0)(1.7,-4.35)(2.6,3.6)

\end{picture}
\end{center}


\noindent\fbox{\parbox{\textwidth}{\rm 
Although this second solution has labeled axes, the graph has no title
so it is not clear to someone else what they are viewing.  There are
also no labels on the extreme points or intercepts and no scale is given
on the $y$-axis.  Although this is a bad example of a solution to this
question it  would be  sufficient if only a rough sketch
was required.
}}
\end{badexample}


Choose a suitable scale so that all of the interesting features of the
graph are included.  For example, when graphing our function
\begin{displaymath}
  y= x^3 - 4x
\end{displaymath}
a suitable scale meant including all of the extreme points and $x$- and
$y$-intercepts.  
Be especially careful about the scale when using a graphing calculator or
computer. They do not have any idea which points are interesting, so they
will often give inappropriate scales.
The following graphs
give examples of bad choices of scale for our problem. 

\begin{badexample}


\setlength{\unitlength}{0.55cm}

\begin{center}
\begin{picture}(10,8)(-5,-4)
% make axes
\put(0,-4){\vector(0,1){8}}
\put(-5,0){\vector(1,0){10}}
\thinlines
%\multiput(-4,-0.1)(1,0){9}{\line(0,1){0.2}}
%\multiput(-0.1,-3.5)(0,0.5){15}{\line(1,0){0.2}}
\put(5.2,0.2){$x$}
\put(0.2,3.8){$y$}

\put(-4.6,-0.6){$-40$}
\put(3.9,-0.6){$40$}
\put(-2,3){$60000$}
\put(-2.5,-4){$-60000$}

\qbezier(-5,-4)(-3,0)(0,0)
\qbezier(5,4)(3,0)(0,0)

\end{picture}
%%%%%%
\hfill
%%%%%%
\begin{picture}(10,8)(-5,-4)

\qbezier(-5,4)(-3,3.5)(0,0)
\qbezier(5,-4)(3,-3.5)(0,0)

% make axes

\put(0,-4){\vector(0,1){8}}
\put(-5,0){\vector(1,0){10}}
\thinlines

\put(5.2,0.2){$x$}
\put(0.2,3.8){$y$}

\put(-4.4,-0.6){$-1$}
\put(3.9,-0.6){$1$}

\put(-1,3){$1$}
\put(-1.5,-4){$-1$}

\end{picture}

\end{center}


\noindent\fbox{\parbox{\textwidth}{\rm 
The graph on the left uses
too large a scale so the extreme points are hard to see and will be
difficult to label.  The scale used for the graph on the right is too 
small as it fails to include two of the $x$-intercepts.
}}
\end{badexample}

Also make sure that what you are graphing makes sense.  If you are drawing
the graph of a function that represents the distance of a space shuttle
from the surface of the earth against time, it makes no sense to include
negative distances, as this is unphysical.  Similar reasoning applies to
graphs representing populations, areas, and volumes.

Finally, being able to draw and sketch graphs well can sometimes help
while trying to answer a question, even if a graph is not required in the
solution.  For more difficult problems in particular, being able to sketch
a function to identify its behavior often provides more insight than
looking at the equation alone.

%
% $Log: introduction.tex,v $
% Revision 1.4  2002/10/21 15:26:44  mjm
% to ohio, revised
%
% Revision 1.3  2001/08/15 18:07:03  mjm
% APPM preprint 466
%
% Revision 1.2  1999/08/09 19:55:27  mjm
% beta release
%
% Revision 1.1  1999/06/10 23:31:49  mjm
% Initial revision
% 
%

%\begin{probdescript}{Writing an Introduction to Your Problem}
\handout{Introductions and Conclusions}

\subsection*{Introductions}

The purpose of  an introduction is to explain the important
techniques used to solve your problem. The length may vary from a few
sentences to a paragraph. The length usually depends upon the number
of steps in your problem needed to get to a solution. A well
constructed introduction will make your answer and work easy to
follow. On the other hand, a poorly constructed introduction will
confuse the reader.  


We will illustrate some things to cover in your introduction with the
problem:
{\em
Use the Sandwich Theorem to find the
      asymptotes of the curve $y= \sin(x)/x$.} 


\begin{itemize}
\item  What the problem is about.
 Ask yourself: ``What am I doing?''
  Someone reading your problem should get an idea of what you are
  about to do. If there are multiple parts then explain each part.

\begin{badexample}
I'm going to calculate the asymptotes.
\end{badexample}

\begin{goodexample}
We will find
  the asymptotes of the function $y=\sin(x)/x$.\\
 \fbox{\rm This one supplies more information to the reader.} 
\end{goodexample}
  
\item The technique(s) you are going to use.
 This will usually correspond to
  the topic you are learning.  Your assigned problem will generally
  come from a section you are covering in class, and you can use this
as a guide.   Sometimes your
  problem will involve several concepts and you should mention what
  the predominant technique is that you are going to use. Do not
  mention the details of your calculations.

\begin{badexample}  
Calculating the limit of $y=\sin(x)/x$ will tell me where
  the asymptotes are.  
\end{badexample}

\begin{goodexample}
To find the asymptotes, we will take the limits as $x$ approaches
infinity or negative infinity. To determine these limits, we will use
the Sandwich Theorem.\\
\fbox{\rm The techniques employed are mentioned}
\end{goodexample}

\item Physical interpretation, if appropriate.
  Whether this
  is necessary depends upon the question. If applicable, this step is
  both helpful for you and the person who is reading your problem.
  Again, think about what you are doing. Give a physical description
  of the mathematical steps you are employing to solve your problem.
  If you were to draw a tangent line at a point, what would the
  picture look like?

\begin{badexample}  
By looking at the graph, $y$ goes to the $x$ axis.
\fbox{\rm What are you talking about?}
\end{badexample}

\begin{goodexample}
By looking at the graph of the function, we see that as $x$ approaches
$\pm\infty$, $y$ approaches~0, so the graph approaches the $x$ axis.\\
\fbox{\rm This description is much clearer.}
\end{goodexample}
 

\end{itemize}
Putting it
  all together: 
%\end{probdescript}

%\clearpage
%\secondpage
%\begin{probexamples}


      \begin{badexample}
        I'm going to calculate the asymptotes. Calculating the limit
        of $y=\sin(x)/x$ will tell me where the asymptotes are. By
        looking at the graph, y goes to the $x$ axis.
      \end{badexample}

      \begin{goodexample} We will find the asymptotes of the function
      $y=\sin(x)/x$.  To find the asymptotes, we will take the limits
      as $x$ approaches infinity or negative infinity. To determine
      these limits, we will use the Sandwich Theorem.  By looking at
      the graph of the function, we see that as $x$ approaches
      $\pm\infty$, $y$ approaches~0, so the graph approaches the $x$
      axis.  \end{goodexample}

\noindent\fbox{\parbox{\textwidth}{\rm 
        The bad introduction doesn't flow very well.  Someone
        reading it would not understand what you did.  It is obvious
        that you calculated the asymptotes, but when solving a problem,
        the techniques employed are as important as your answer. The
        good introduction briefly explains how you used the Sandwich
        Theorem along with the result you obtained. This would allow the
        reader to have a clear understanding of the mathematical
        calculations that follow.
}}

%\secondpage

%\subsection*{Examples:}

%{\bf Problem:}
%	{\em  Find $dy/dx$ if 
%      \begin{displaymath}
%        y=\int_1^{x^2} \cos(t)dt.
%      \end{displaymath}
%	}

%      \begin{badexample}
%        I'm going to take the derivative of this integral to find
%        $dy/dx$. Using the Chain Rule I found $dy/dx = 2x \cos(x^2).$
%    \end{badexample}

%\bigskip

%    \begin{goodexample}
%       I'm going to use the Fundamental Theorem of Calculus to
%      find $dy/dx$. By treating $y$ as a composite of 
%	$y=\int_1^u \cos(t)d$  and
%      $u=x^2$, I can use the Chain Rule to find $dy/dx = dy/du \times
%      du/dx$.  Evaluating these results in the answer $dy/dx = 2x
%      \cos(x^2)$.
%    \end{goodexample}

%\bigskip

%\noindent\fbox{\parbox{\textwidth}{\rm 
%      The bad example would not give the reader a clear guide to
%      follow your calculations.  In the good example, the
%      Fundamental theorem of Calculus is mentioned along with a
%      comment on how $y$ is going to be treated.  Thus when you do your
%      calculations, each step through the Chain Rule will have a
%      logical flow.
%}}


%\bigskip
%{\bf Problem:}
%	{\em 
%      The cost function at
%      XYZ Co. is $c(x)=x^3-6x^2+15x$. Is there a production level that
%      minimizes average cost? If so, what is it?  
%}

%\bigskip

%      \begin{badexample}
%        I'm going to calculate the derivative of the function to
%        find the marginal cost function. When this function equals the
%        average cost function we have a production level which
%        minimizes average cost.
%    \end{badexample}

%\bigskip

%    \begin{goodexample}
%      The marginal cost function is calculated as the first derivative
%      of the cost function.  The average cost function is $c(x)/x$. By
%      setting these two equations equal to each other and solving for
%      $x$, I get the production levels that might minimize the average
%      cost function. By performing the first and second derivative
%      tests on the average cost function at $x$ I can verify if the
%      production level $x$ is in fact a minimum.
%\marginpar{funny answer}
%    \end{goodexample}      

%\bigskip

%\noindent\fbox{\parbox{\textwidth}{\rm 
%        Clearly, this introduction provides a better outline about the
%        techniques used to solve the problem.
%}}


%\marginpar{conclusion}

\secondpage
\subsection*{Conclusions}

Your solution should end with a brief concluding paragraph. The
conclusion is a short, concise paragraph that summarizes the main
ideas, techniques and results discussed in the problem.  The
concluding paragraph should focus on the main results and main ideas
so that the reader can absorb the important parts of your solution to
the problem. The concluding paragraph should not
introduce any new material but simply bring back the reader to the main ideas
of your solution to the problem. This means that you should not
introduce something in the conclusion that was not discussed earlier in
your solution.

When you write a conclusion, put yourself in the potential reader's
position and imagine what kind of information this reader may look for
when he or she reads the concluding paragraph.  Also, ask yourself
``What do I want my reader to remember?''  Here are a few suggestions
of what you may want to include in the conclusion paragraph:

\begin{itemize}
\item Restate the problem you solved.
\item Restate the results and interpretation of the results. If your
  results are lengthy, such as a big table with numerical values,
  just  describe the results qualitatively. For example,
  give the range of values for your results.
\item Indicate if your results seem reasonable. If not, then explain why.
\item If you were unsuccessful in solving your problem, mention
  possible reasons for this.
\item If applicable, mention other problems that are related to the
  problem you solved.
\item Give suggestions for improvements, i.e., what else you could do
  for the problem.
\end{itemize}

These are just suggestions, and not all of them may be applicable for
your problem. 


%\secondpage

%\subsection*{Examples:}

{\bf Example Problem:}
  \emph{Find
        the dimensions of a right circular cylinder of volume
        $10^{-6}\ m^3$ such that a minimum amount of material is
        used.}  

\bigskip

      \begin{badexample}
        In this problem I used calculus to solve a mathematical
        problem.
\fbox{\rm This gives no information.}
    \end{badexample}

\bigskip

    \begin{badexample}
      I solved for $h$ in $\pi r^2h=10^{-6}$ and plugged the
      expression for $h$ into $2\pi r^2+2\pi rh$. Then I took the
      derivative of $2\pi r^2+\frac{2\times 10^{-6}}{r}$. Then I set
      the derivative equal to zero. Then I solved for $r$.
    \end{badexample}


\noindent\fbox{\parbox{\textwidth}{\rm 
      The last example fails to focus on the main ideas of the
      problem. It gives a (bad) summary of some of the steps used to
      solve the problem. It does not summarize the problem given to us
      and it does not clearly indicate what your bottom line was or
      if your were even successful in solving the problem.
}}

\bigskip

    \begin{goodexample}
      The problem discussed in this paper is to find the dimensions of
      a right circular cylinder of volume $10^{-6}\ m^3$ such that a
      minimum amount of material is used. The problem was solved by
      minimizing the area of the cylinder using the first and second
      derivative test. The optimal dimensions were found to be radius
      $\cong 5.42\ cm$ and height $\cong 10.84\ cm$. The method used
      for solving the problem can be applied to other optimization
      problems, e.g., finding the dimensions of other shapes with a
      fixed volume such that the used material is minimized. A
      possible development of the technique used could be to consider
      optimization of a function of more than two parameters and more
      than one constraint.
    \end{goodexample}

\noindent\fbox{\parbox{\textwidth}{\rm 
      Note that the suggested conclusion requires that we actually did
      discuss possible applications and developments of the technique
      earlier in the report, something that may not be required for
      you in your course.
}}

%\marginpar{conclusion}
%%
% $Log: conclusion.tex,v $
% Revision 1.3  2002/01/08 18:52:29  mjm
% that
%
% Revision 1.2  1999/07/28 17:23:11  mjm
% version sent back to Kristian
%
% Revision 1.1  1999/06/10 23:33:32  mjm
% Initial revision
%
%

\handout{Conclusions}

Your solution should end with a brief concluding paragraph. The
conclusion is a short, concise paragraph that summarizes the main
ideas, techniques and results discussed in the problem.  The
concluding paragraph should focus on the main results and main ideas
so that the reader can absorb the important parts of your solution to
the problem. The concluding paragraph should not
introduce any new material but simply bring back the reader to the main ideas
of your solution to the problem. This means that you should not
discuss something in the conclusion that was not discussed earlier in
your solution.

When you write a conclusion, put yourself in the potential reader's
position and imagine what kind of information this reader may look for
when he or she reads the concluding paragraph.  Also, ask yourself
``What do I want my reader to remember?''  Here are a few suggestions
of what you may want to include in the conclusion paragraph:

\begin{itemize}
\item Restate the problem you solved.
\item Restate the results and interpretation of the results. If your
  results are lengthy, such as a big table with numerical values,
  just  describe the results qualitatively. For example,
  give the range of values for your results.
\item Indicate if your results seem reasonable. If not, then explain why.
\item If you were unsuccessful in solving your problem, mention
  possible reasons for this.
\item If applicable, mention other problems that are related to the
  problem you solved.
\item Give suggestions for improvements, i.e., what else you could do
  for the problem.
\end{itemize}

These are just suggestions, and not all of them may be applicable for
your problem.  As mentioned above, the conclusion should be short and
concise without discussing anything new that was not discussed earlier
in the report.


\secondpage

\subsection*{Examples:}

{\bf Problem:}
  \emph{Find
        the dimensions of a right circular cylinder of volume
        $10^{-6}\ m^3$ such that a minimum amount of material is
        used.}  

\bigskip

      \begin{badexample}
        In this problem I used calculus to solve a mathematical
        problem.\\
\fbox{This gives no information.}
    \end{badexample}

\bigskip

    \begin{badexample}
      I solved for $h$ in $\pi r^2h=10^{-6}$ and plugged the
      expression for $h$ into $2\pi r^2+2\pi rh$. Then I took the
      derivative of $2\pi r^2+\frac{2\times 10^{-6}}{r}$. Then I set
      the derivative equal to zero. Then I solved for $r$.
    \end{badexample}


\noindent\fbox{\parbox{\textwidth}{\rm 
      The last example fails to focus on the main ideas of the
      problem. It gives a (bad) summary of some of the steps used to
      solve the problem. It does not summarize the problem given to us
      and it does not clearly indicates what your bottom line was or
      if your were even successful in solving the problem.
}}

\bigskip

    \begin{goodexample}
      The problem discussed in this paper is to find the dimensions of
      a right circular cylinder of volume $10^{-6}\ m^3$ such that a
      minimum amount of material is used. The problem was solved by
      minimizing the area of the cylinder using the first and second
      derivative test. The optimal dimensions were found to be radius
      $\cong 5.42\ cm$ and height $\cong 10.84\ cm$. The method used
      for solving the problem can be applied to other optimization
      problems, e.g., finding the dimensions of other shapes with a
      fixed volume such that the used material is minimized. A
      possible development of the technique used could be to consider
      optimization of a function of more than two parameters and more
      than one constraint.
    \end{goodexample}

\noindent\fbox{\parbox{\textwidth}{\rm 
      Note that the suggested conclusion requires that we actually did
      discuss possible applications and developments of the technique
      earlier in the report, something that may not be required for
      you in your course.
}}

\marginpar{conclusion}

\end{document}
%%%%%%%%%%%%%%%%%%%%%%%%%%%%%%%%%%%%%%%%%%%%%%%%%%%%%%%%%%%%%%%%%%%%%%%%%
%%%%%%%%%%%%%%%%%%%%%%%%%%%%%%%%%%%%%%%%%%%%%%%%%%%%%%%%%%%%%%%%%%%%%%%%%
%%%%%%%%%%%%%%%%%%%%%%%%%%%%%%%%%%%%%%%%%%%%%%%%%%%%%%%%%%%%%%%%%%%%%%%%%
